\documentclass{article}%
\usepackage[T1]{fontenc}%
\usepackage[utf8]{inputenc}%
\usepackage{lmodern}%
\usepackage{textcomp}%
\usepackage{lastpage}%
\usepackage{geometry}%
\geometry{margin=2cm}%
\usepackage{expex}%
\usepackage[T1]{fontenc}%
\usepackage[utf8]{inputenc}%
\usepackage[spanish]{babel}%
%
\title{Análisis Lingüístico: Gramatica Normativa Kaqchikel}%
\author{Generado por el TFM de Elías Carrasco}%
%
\begin{document}%
\normalsize%
\maketitle%
\section{Page 1}%
\label{sec:Page1}%
\centerline{\textbf{Nuk’b’il Yol Mam}}%
K’ulb’il Yol Mam te K’ulb’il Yol Twitz Paxil 230 %
\subsection{Poq’b’il t{-}xilen yol / Extensiones semánticas}%
\label{subsec:Poqbilt{-}xilenyol/Extensionessemnticas}%

%
Los cambios semánticos se toman cuando una palabra ha tomado otros significados distintos a los originales o comunes. %
Algunas palabras son similares en pronunciación pero vienen de diferentes fuentes. %
Las razones más comunes del cambio semántico provienen de la pronunciación y/o en extensiones de variantes específicos en el contexto que se habla. %
Los cambios semánticos se dan sobre todo en sustantivos, adjetivos y verbos. %
\subsection{B’ib’aj/Sustantivos}%
\label{subsec:Bibaj/Sustantivos}%

%
Los cambios ocurren en las formas absolutas de las palabras como en su forma poseída, la mayoría de los cambios se deben a procesos fonológicos. %
Ejemplos: %
\par\medskip%
\ex
\textit{tu sombrero} \\
\begingl
\gla Tpasa //
\glb  //
\glc  //
\endgl
\xe
%
\medskip%
\par\medskip%
\ex
\textit{tu faja} \\
\begingl
\gla Tpasa //
\glb  //
\glc  //
\endgl
\xe
%
\medskip%
\par\medskip%
\ex
\textit{estrella} \\
\begingl
\gla Che’w //
\glb  //
\glc  //
\endgl
\xe
%
\medskip%
\par\medskip%
\ex
\textit{frío} \\
\begingl
\gla Che’w //
\glb  //
\glc  //
\endgl
\xe
%
\medskip%
\par\medskip%
\ex
\textit{piedra} \\
\begingl
\gla Xaq //
\glb  //
\glc  //
\endgl
\xe
%
\medskip%
\par\medskip%
\ex
\textit{barranco} \\
\begingl
\gla Xaq //
\glb  //
\glc  //
\endgl
\xe
%
\medskip%
\par\medskip%
\ex
\textit{piel} \\
\begingl
\gla Chi’yaj //
\glb  //
\glc  //
\endgl
\xe
%
\medskip%
\par\medskip%
\ex
\textit{carnoso} \\
\begingl
\gla Chi’yaj //
\glb  //
\glc  //
\endgl
\xe
%
\medskip

%
\end{document}